\section{Two obstructions to the direct tensor repairs of F-SM34}

\paragraph{Scope.}
These are obstructions to specified constructions, not a proof that
the additive category in F-SM34 admits no symmetric monoidal
structure.  In particular, they do not turn that additive
counterexample into a counterexample to all the handwritten
conditions of cote~134-8.  The full tensor, its scalar unit, and the
interval remain separate requirements.

\subsection{F-SM36: no normalized contraction of the double-dual quotient}

Let \(k\) be a field, let \(V^*=\operatorname{Hom}_k(V,k)\), and
put \(E(V)=V^{**}/\operatorname{ev}_V(V)\).  There is no family of
linear maps
\[
 c_{V,W}:E(V\otimes_k W)\longrightarrow E(V)\otimes_k W
                                                               \tag{SM36.1}
\]
natural in both vector spaces and equal to the identity for
\(W=k\), under the usual unit identifications.  In fact, for any
fixed \(V\) with \(E(V)\ne0\), naturality in \(W\) alone is
already impossible.

\begin{proof}
First such a \(V\) exists.  Take \(V=k^{(\mathbb N)}\), so that
\(V^*=k^{\mathbb N}\).  The quotient
\(k^{\mathbb N}/k^{(\mathbb N)}\) is nonzero, as witnessed by
the constant sequence \(1\).  Choose a nonzero linear functional
on this quotient and pull it back to \(f\in V^{**}\).
If \(f\) were evaluation at \(v\in V\), its vanishing on every
finite-support coordinate vector of \(V^*\) would force every
coordinate of \(v\) to vanish.  This would give \(f=0\), a
contradiction.  Thus \(e=[f]\in E(V)\) is nonzero.

The rest of the argument applies to any \(V\) and any
\(f\in V^{**}\) with nonzero class \(e\).  Set
\(W=k^{(\mathbb N)}\), with basis \((w_n)\), and identify
\[
        (V\otimes W)^*\simeq\prod_{n\ge0}V^*.
\]
On the subspace \(\bigoplus_{n\ge0}V^*\) of this product,
define the functional
\[
             (\lambda_n)\longmapsto\sum_n f(\lambda_n).
\]
The sum is finite on this subspace.  Extend this functional
linearly to the whole product, obtaining
\(F\in(V\otimes W)^{**}\), and write \(u=[F]\).
For the coordinate functional \(\epsilon_n:W\to k\), the
dual of \(1_V\otimes\epsilon_n\) is the inclusion of the
\(n\)-th copy of \(V^*\) into \(\prod_mV^*\).  Consequently
\[
       E(1_V\otimes\epsilon_n)(u)=e\qquad(n\ge0).             \tag{SM36.2}
\]
Naturality and the stated normalization of \(c\) would give
\[
      (1_{E(V)}\otimes\epsilon_n)c_{V,W}(u)=e
                      \qquad(n\ge0).                        \tag{SM36.3}
\]
Every element of \(E(V)\otimes k^{(\mathbb N)}\), however,
has only finitely many nonzero coordinates.  Equations
(SM36.3) contradict \(e\ne0\).  The linear extensions used
here are ordinary algebraic extensions of linear functionals;
no topology or continuity is involved.
\end{proof}

\paragraph{Consequence for the actual additive category.}
Write \(\mathcal N\) for F-SM34's category of
\[
 (X_0,X_1,b:X_0\to X_1,a:E(X_1)\to X_0),
 \qquad ba=0,\quad aE(b)=0,
\]
and \(\kappa(W)=(W,0)\).  There is no unital bifunctor
\(\star:\mathcal N\times\mathcal N\to\mathcal N\) with
unit \(\kappa(k)\) which has natural identifications
\[
 U_i(X\star\kappa(W))\simeq U_i(X)\otimes_kW
                  \quad(i=0,1),                             \tag{SM36.4}
\]
compatible with the canonical vector-space unit identifications.
Here \(U_i\) extracts the indicated component.  This conclusion
does not even require associativity or symmetry of \(\star\).

\begin{proof}
For every \(V\) the object
\(P_1(V)=(E(V),V,0,\mathrm{id})\) belongs to \(\mathcal N\),
and depends functorially on \(V\).  Under (SM36.4), the
\(a\)-operation of \(P_1(V)\star\kappa(W)\) is a map
exactly of the form (SM36.1).  Functoriality makes it natural
in \(V,W\).  Compatibility of the unit isomorphism, which is
an isomorphism in \(\mathcal N\), makes its value at \(W=k\)
the identity on \(E(V)\).  F-SM36 rules this out.
\end{proof}

This leaves open a tensor whose underlying vector spaces are
enlarged.  In particular, (SM36.4) is not a consequence of the
mere bilinearity written in the original scan.  A tensor
preserving scalar copowers in \(\mathcal N\) would involve
the actual coproducts of \(\mathcal N\), whose components
are not generally ordinary direct sums.

\subsection{F-SM37: the arrow monad has no usual tensorial strength}

Let \(\mathcal A=\operatorname{Fun}([0\to1],\operatorname{Vect}_k)\).
Regard \(C=(C_0\xrightarrow{b}C_1)\) as a cochain complex
in degrees \(0,1\), and give \(\mathcal A\) its truncated
cochain tensor
\[
 \begin{aligned}
 (C\odot D)_0&=C_0\otimes D_0,\\
 (C\odot D)_1&=(C_1\otimes D_0)\oplus(C_0\otimes D_1),\\
 b_{C\odot D}&=(b_C\otimes1,\ 1\otimes b_D).
 \end{aligned}                                               \tag{SM37.1}
\]
Its unit is \(I=(k\to0)\).  Associativity and symmetry are
the restrictions of the cochain tensor: in an arbitrary finite
product the degree-zero term has only degree-zero factors, and
the degree-one term has precisely one degree-one factor.

For any nonzero additive exact endofunctor \(E\) of vector
spaces, define
\[
 H(C)=(E(\operatorname{coker}b_C)\to0),\qquad
 T=\mathrm{id}_{\mathcal A}\oplus H.                         \tag{SM37.2}
\]
Then \(H^2=0\), so \(T\) is the square-zero monad whose unit
is the first inclusion and whose multiplication adds the two
\(H\)-summands.  For \(E(V)=V^{**}/V\), its algebra category
is precisely \(\mathcal N\) above.  Nevertheless there is no
natural family
\[
       s_{C,D}:T(C)\odot D\longrightarrow T(C\odot D)
                                                               \tag{SM37.3}
\]
which is the identity when \(D=I\).  Thus \(T\) cannot be
a strong monad for (SM37.1); in particular this model cannot
obtain a tensor by invoking a commutative strong monad on that
base category.

\begin{proof}
Since the cokernel of \(E(\operatorname{coker}b)\to0\) is
zero, \(H^2=0\).  The monad identities follow by adding the
three copies of \(H(C)\) in \(T^3(C)\); all composites of two
\(H\)'s vanish.  An algebra structure has the form
\(1_C\oplus\bar a:C\oplus H(C)\to C\).  The requirement
that it be an arrow morphism is \(b_C\bar a=0\), and the
monad multiplication condition is automatic because \(H^2=0\).
Exactness of \(E\) identifies a map
\(E(\operatorname{coker}b_C)\to C_0\) with a map
\(a:E(C_1)\to C_0\) satisfying \(aE(b_C)=0\).
This gives the stated algebra-category identification, including
its morphisms.

Choose \(V\) with \(E(V)\ne0\), and write
\[
 S=(0\to V),\qquad D(V)=(V\xrightarrow{1}V),\qquad
 D=D(k).
\]
Let \(j:S\to D(V)\) be zero in degree zero and the identity
in degree one.  Let \(q:D\to I\) be the identity in degree
zero and zero in degree one.  Directly from (SM37.1),
\[
 \begin{aligned}
 S\odot D&=S, & S\odot q&=1_S,\\
 T(S)&=(E(V)\xrightarrow{0}V), & T(D(V))&=D(V).
 \end{aligned}
                                                               \tag{SM37.4}
\]
Naturality of \(s\) with respect to \(q\), and the
normalization at \(I\), force
\[
                    (s_{S,D})_0=1_{E(V)}.                    \tag{SM37.5}
\]
On the other hand,
\[
 D(V)\odot D=(V\xrightarrow{v\mapsto(v,v)}V\oplus V).
\]
The map \(j\odot1_D\) sends the degree-one term \(V\)
to the first summand of \(V\oplus V\).  Its induced map
on cokernels is the isomorphism
\[
       V\longrightarrow(V\oplus V)/\operatorname{diag}(V).
\]
It follows that the degree-zero map of \(T(j\odot1_D)\)
is an isomorphism from \(E(V)\) onto the \(H\)-summand of
\(T(D(V)\odot D)_0\), and hence is nonzero.
But \(T(j)_0:E(V)\to V\) is zero.  The naturality square
for \(j\) and \(D\) would therefore give
\[
 T(j\odot1_D)_0(s_{S,D})_0
   =(s_{D(V),D})_0(T(j)\odot1_D)_0=0,
\]
contradicting (SM37.5).  No associativity axiom for a strength
was used: naturality and its unit axiom already contradict
each other.
\end{proof}

\paragraph{What is still open.}
F-SM36 excludes an ordinary componentwise scalar tensor, and
F-SM37 excludes the indicated strong-monad construction.
Neither excludes a different tensor on \(\mathcal N\), a
different monoidal base for an infinitary theory, or another
category satisfying all the original assumptions.  Conversely,
neither proves that the original interval forces the missing
derived normalization on infinite scalar modules.  These are
the exact remaining alternatives.

\paragraph{Bounded literature check on a different proposed route.}
An ordinary Hopf-algebra structure does not automatically supply
the needed tensor on a contramodule category.  Johnston (2026,
Section~2.3) explicitly constructs an action involving a comodule
and a contramodule, using \(\operatorname{Hom}_k\), rather than
a tensor of two arbitrary contramodules.  Consequently a
completed Hopf-algebra candidate still needs its own construction
of the symmetric tensor, its trivial-module unit, and both
scalar adjoints.  This observation is not a nonexistence theorem
for all such completed constructions.

\paragraph{Reference.}
Johnston, D. (2026). Contramodules for algebraic groups: The
existence of mock projectives. \emph{Canadian Mathematical
Bulletin, 69}(1), 184--200. Published online August~4, 2025.
\url{https://doi.org/10.4153/S0008439525100933}.
Section~2.3, introductory paragraph and displayed
\(\operatorname{Hom}_k\) bifunctor.  The elementary proofs
F-SM36--37 above do not use that paper.

\paragraph{Two further repairs tested; no new original-status conclusion.}
The intake's three-sort proposal refers to an unavailable full
\texttt{C.tex} and an unrun coherence script.  For the radical-square-zero
functor \(H(U_0,U_1,U_2)=(E(U_2),U_0,U_1)\), its free projectives are
\(P_0(V)=(V,V,0)\), \(P_1(V)=(0,V,V)\), and
\(P_2(V)=(EV,0,V)\), with the corresponding identity structure maps.
Merely substituting these into the standard truncated-cochain tensor
does not work.  For a nonzero ghost
\(a:P_0(k)\to P_2(W)\) and the map
\(b:P_2(W)\to P_1(W)\) that is the identity in the last component,
one has \(ba=0\).  Tensoring with \(P_0(k)\) under the proposed
standard disk rules gives \((a,0)\) followed by
\(\binom{b}{\pm1}\).  Their composite has nonzero second component
\(\pm a\).  This excludes only these specific rules, not the
unavailable intake construction.

There is a separate defect in the one-sort square-zero repair.
For an exact linear functor \(E\) on vector spaces, the monad
\(T=\mathrm{id}\oplus E\) with zero \(E^2\)-multiplication has
the usual commutative square-zero strengths.  An algebra is a map
\(a:EV\to V\) with \(aE(a)=0\).  Suppose \(E\) vanishes on
finite-dimensional spaces and is nonzero on some \(V\).
The trivial algebra \(\kappa(k)\) is free, since \(Ek=0\).
If the trivial-algebra embedding \(\kappa\) had a right adjoint,
it would preserve coproducts; its value on every free vector space
\(V\) would therefore have to be the free algebra \(T(V)\).
But \(T(V)\) has nonzero structure map when \(EV\ne0\), while
\(\kappa(V)\) has zero structure map.  Thus this repair lacks
one of the required scalar adjoints.

Its usual algebra tensor also fails scalar compatibility.  The
tensor of two trivial algebras is the quotient of the free algebra
on \(V\otimes W\) by the images of the two strengths; explicitly
its underlying vector space is
\[
 (V\otimes W)\oplus E(V\otimes W)/R,\qquad
 R=\operatorname{im}(EV\otimes W)
       +\operatorname{im}(V\otimes EW).
\]
The structure map on the second summand is the quotient map from
\(E(V\otimes W)\); the square-zero condition adds no further
relation there.  For \(W=V\), choose a basis \((v_i)\) and define
\(\delta(v_i)=v_i\otimes v_i\) and
\(m(v_i\otimes v_j)=\delta_{ij}v_i\).  Then \(m\delta=1\).
For each fixed \(v\in V\), the maps \(m(-\otimes v)\) and
\(m(v\otimes-)\) have finite-dimensional image.  Thus \(E(m)\)
kills both strength images, whereas
\(E(m)E(\delta)=1_{EV}\).  Consequently
\(E(V\otimes V)/R\ne0\) if \(EV\ne0\), and the algebra tensor
cannot be \(\kappa(V\otimes V)\).  These two failed repairs do
not justify adding tensor exactness, closedness or AB4 to the
original finite-additivity assumptions.
